长江流域是我国淡水渔业最发达、生物多样性最丰富的区域之一，但过度捕捞、工业废水排放、水域塑料污染长期存在，葛洲坝、三峡大坝等水利工程还造成鱼类洄游通道阻断。受这些因素共同影响，渔业资源持续衰退，典型食物链结构逐渐退化。自2021年1月1日实施长江十年禁渔政策后，局部水域水质开始改善，水草等水生植被逐步恢复。2025年监测数据显示，青、草、鲢、鳙四大家鱼资源量已恢复至禁渔前的约1.8倍，这一结果表明禁渔对基础资源和经济鱼类具有显著的修复作用。
